\documentclass[10pt,a4paper]{amsart}
\usepackage{amsmath,amssymb,amsthm,graphicx}
% For commutative diagrams we use Xy-pic
\usepackage[all]{xy}

% We use 2cell for 2-commutative diagrams.
\xyoption{2cell}
% Theorem environments.
%
\theoremstyle{plain}
\newtheorem{theorem}[subsection]{Theorem}
\newtheorem{proposition}[subsection]{Proposition}
\newtheorem{lemma}[subsection]{Lemma}

\theoremstyle{definition}
\newtheorem{definition}[subsection]{Definition}
\newtheorem{example}[subsection]{Example}
\newtheorem{exercise}[subsection]{Exercise}
\newtheorem{situation}[subsection]{Situation}

\theoremstyle{remark}
\newtheorem{remark}[subsection]{Remark}
\newtheorem{remarks}[subsection]{Remarks}

\numberwithin{equation}{subsection}

% Macros
%
\def\lim{\mathop{\mathrm{lim}}\nolimits}
\def\colim{\mathop{\mathrm{colim}}\nolimits}
\def\Spec{\mathop{\mathrm{Spec}}}
\def\Hom{\mathop{\mathrm{Hom}}\nolimits}
\def\Ext{\mathop{\mathrm{Ext}}\nolimits}
\def\SheafHom{\mathop{\mathcal{H}\!\mathit{om}}\nolimits}
\def\SheafExt{\mathop{\mathcal{E}\!\mathit{xt}}\nolimits}
\def\Sch{\mathit{Sch}}
\def\Mor{\mathop{\mathrm{Mor}}\nolimits}
\def\Ob{\mathop{\mathrm{Ob}}\nolimits}
\def\Sh{\mathop{\mathit{Sh}}\nolimits}
\def\NL{\mathop{N\!L}\nolimits}
\def\CH{\mathop{\mathrm{CH}}\nolimits}
\def\proetale{{pro\text{-}\acute{e}tale}}
\def\etale{{\acute{e}tale}}
\def\QCoh{\mathit{QCoh}}
\def\Ker{\mathop{\mathrm{Ker}}}
\def\Im{\mathop{\mathrm{Im}}}
\def\Coker{\mathop{\mathrm{Coker}}}
\def\Coim{\mathop{\mathrm{Coim}}}

% Boxtimes
%
\DeclareMathSymbol{\boxtimes}{\mathbin}{AMSa}{"02}

%
% Macros for moduli stacks/spaces
%
\def\QCohstack{\mathcal{QC}\!\mathit{oh}}
\def\Cohstack{\mathcal{C}\!\mathit{oh}}
\def\Spacesstack{\mathcal{S}\!\mathit{paces}}
\def\Quotfunctor{\mathrm{Quot}}
\def\Hilbfunctor{\mathrm{Hilb}}
\def\Curvesstack{\mathcal{C}\!\mathit{urves}}
\def\Polarizedstack{\mathcal{P}\!\mathit{olarized}}
\def\Complexesstack{\mathcal{C}\!\mathit{omplexes}}
% \Pic is the operator that assigns to X its picard group, usage \Pic(X)
% \Picardstack_{X/B} denotes the Picard stack of X over B
% \Picardfunctor_{X/B} denotes the Picard functor of X over B
\def\Pic{\mathop{\mathrm{Pic}}\nolimits}
\def\Picardstack{\mathcal{P}\!\mathit{ic}}
\def\Picardfunctor{\mathrm{Pic}}
\def\Deformationcategory{\mathcal{D}\!\mathit{ef}}
\usepackage{cleveref}
\crefname{originalremark}{original remark}{original remarks}
\Crefname{originalremark}{Original remark}{Original remarks}
\makeatletter
\newlabel{remark-more-elementary}{{02PK}{source}{}{sourcepointer-1}{}}
\newlabel{remark-more-elementary@cref}{{[originalremark][0][]02PK}{[1][0][]source}{}{}{}}
\makeatother
\title{Multiplicativity of the determinant of a (2,1)-periodic complex under the factorization of one of the three endomorphisms}
\author{The Stacks Project Authors}
\date{stacks-project:chow.tex:local-snapshot-20261012; GNU FDL 1.2 or later}
\begin{document}
\maketitle
\noindent\textit{Editorial note.} All delivered mathematical text is the author's own, copied from the pinned Stacks Project chapter file chow.tex; no statement or proof step was rewritten, reordered or removed.\par\medskip
\noindent\textit{Editorial note.} Preamble provenance: the pinned chapter file opens with an input of the shared preamble and carries no preamble text of its own, so the project's single Stacks-wide preamble staged at sources/stacks-chow/preamble.tex is excerpted verbatim as two spans, lines 14-18 (Xy-pic setup) and lines 154-217 (the native theorem declarations and mathematical macros). Nothing was added, removed or redefined; the wrapper class and the amsmath/amssymb/amsthm/graphicx base packages are supplied by the assembler.\par\medskip
\noindent\textit{Editorial note.} The counted statement and proof contain no cross reference, no citation and no included file, so the delivered body is self contained.\par\medskip
\noindent\textit{Editorial note.} Canonical location of the counted proposition: Stacks tag 02PU (sources/stacks-chow/pinned-original-source.tex line 16983).\par\medskip
\noindent\textit{Allowed standard prerequisites (editorial).}\begin{itemize}
\item Determinant of a (2,1)-periodic complex and the characterization of it by a single symbol comparison. Same chapter, chow.tex, the definition of det\_kappa(M, alpha, beta) for an exact (2, 1)-periodic complex and the subsequent characterisation results. Supplies that det\_kappa of an exact (2, 1)-periodic complex is an element of kappa\textasciicircum{}* and that it is uniquely determined by comparing two symbols, which is the form in which the three determinants are compared in the counted proof.
\item The characterisation of the determinant by any admissible generating sequence. Same chapter, Remark labelled remark-more-elementary (source line 16632), invoked at proof lines 17065, 17077 and 17090. Supplies that the defining relation for D, D\_1 and D\_2 may be written with the particular admissible sequences chosen in the proof; the choice of those sequences, the length bookkeeping and the combination of the three relations stay inside the counted span.
\end{itemize}\medskip
\ifdefined\hypertarget\hypertarget{sourcepointer-1}{}\fi
\noindent\textit{Source reference.} Original remark 02PK: sources/stacks-chow/pinned-original-source.tex, source lines 16632--16632. Same-chapter pointer into the author environment remark at source line 16632; invoked at proof line 17065 of the counted proof. It is retained only as a navigation pointer; every load-bearing step is inside the delivered span.\par\medskip
\begin{lemma}
\label{lemma-multiplicativity-determinant}
Let $R$ be a local ring with residue field $\kappa$.
Let $M$ be a finite length $R$-module.
Let $\alpha, \beta, \gamma$ be endomorphisms of $M$.
Assume that
\begin{enumerate}
\item $I_\alpha = K_{\beta\gamma}$, and similarly for any permutation
of $\alpha, \beta, \gamma$,
\item $K_\alpha = I_{\beta\gamma}$, and similarly for any permutation
of $\alpha, \beta, \gamma$.
\end{enumerate}
Then
\begin{enumerate}
\item The triple $(M, \alpha, \beta\gamma)$
is an exact $(2, 1)$-periodic complex.
\item The triple $(I_\gamma, \alpha, \beta)$
is an exact $(2, 1)$-periodic complex.
\item The triple $(M/K_\beta, \alpha, \gamma)$
is an exact $(2, 1)$-periodic complex.
\item We have
$$
\det\nolimits_\kappa(M, \alpha, \beta\gamma)
=
\det\nolimits_\kappa(I_\gamma, \alpha, \beta)
\det\nolimits_\kappa(M/K_\beta, \alpha, \gamma).
$$
\end{enumerate}
\end{lemma}

\begin{proof}
It is clear that the assumptions imply part (1) of the lemma.

\medskip\noindent
To see part (1) note that the assumptions imply that
$I_{\gamma\alpha} = I_{\alpha\gamma}$, and similarly for kernels
and any other pair of morphisms.
Moreover, we see that
$I_{\gamma\beta} =I_{\beta\gamma} = K_\alpha \subset I_\gamma$ and
similarly for any other pair. In particular we get a short exact sequence
$$
0 \to I_{\beta\gamma} \to I_\gamma \xrightarrow{\alpha} I_{\alpha\gamma} \to 0
$$
and similarly we get a short exact sequence
$$
0 \to I_{\alpha\gamma} \to I_\gamma \xrightarrow{\beta} I_{\beta\gamma} \to 0.
$$
This proves $(I_\gamma, \alpha, \beta)$ is an exact $(2, 1)$-periodic
complex. Hence part (2) of the lemma holds.

\medskip\noindent
To see that $\alpha$, $\gamma$ give well defined endomorphisms
of $M/K_\beta$ we have to check that $\alpha(K_\beta) \subset K_\beta$
and $\gamma(K_\beta) \subset K_\beta$. This is true because
$\alpha(K_\beta) = \alpha(I_{\gamma\alpha}) = I_{\alpha\gamma\alpha}
\subset I_{\alpha\gamma} = K_\beta$, and similarly in the other case.
The kernel of the map $\alpha : M/K_\beta \to M/K_\beta$ is
$K_{\beta\alpha}/K_\beta = I_\gamma/K_\beta$. Similarly,
the kernel of $\gamma : M/K_\beta \to M/K_\beta$ is equal to
$I_\alpha/K_\beta$. Hence we conclude that (3) holds.

\medskip\noindent
We introduce $r = \text{length}_R(K_\alpha)$,
$s = \text{length}_R(K_\beta)$ and $t = \text{length}_R(K_\gamma)$.
By the exact sequences above and our hypotheses we have
$\text{length}_R(I_\alpha) = s + t$, $\text{length}_R(I_\beta) = r + t$,
$\text{length}_R(I_\gamma) = r + s$, and
$\text{length}(M) = r + s + t$.
Choose
\begin{enumerate}
\item an admissible sequence $x_1, \ldots, x_r \in K_\alpha$
generating $K_\alpha$
\item an admissible sequence $y_1, \ldots, y_s \in K_\beta$
generating $K_\beta$,
\item an admissible sequence $z_1, \ldots, z_t \in K_\gamma$
generating $K_\gamma$,
\item elements $\tilde x_i \in M$ such that $\beta\gamma\tilde x_i = x_i$,
\item elements $\tilde y_i \in M$ such that $\alpha\gamma\tilde y_i = y_i$,
\item elements $\tilde z_i \in M$ such that $\beta\alpha\tilde z_i = z_i$.
\end{enumerate}
With these choices the sequence
$y_1, \ldots, y_s, \alpha\tilde z_1, \ldots, \alpha\tilde z_t$
is an admissible sequence in $I_\alpha$ generating it.
Hence, by Remark \ref{remark-more-elementary} the determinant
$D = \det_\kappa(M, \alpha, \beta\gamma)$ is the
unique element of $\kappa^*$ such that
\begin{align*}
[y_1, \ldots, y_s,
\alpha\tilde z_1, \ldots, \alpha\tilde z_s,
\tilde x_1, \ldots, \tilde x_r] \\
= (-1)^{r(s + t)} D
[x_1, \ldots, x_r,
\gamma\tilde y_1, \ldots, \gamma\tilde y_s,
\tilde z_1, \ldots, \tilde z_t]
\end{align*}
By the same remark, we see that
$D_1 = \det_\kappa(M/K_\beta, \alpha, \gamma)$
is characterized by
$$
[y_1, \ldots, y_s,
\alpha\tilde z_1, \ldots, \alpha\tilde z_t,
\tilde x_1, \ldots, \tilde x_r]
=
(-1)^{rt} D_1
[y_1, \ldots, y_s,
\gamma\tilde x_1, \ldots, \gamma\tilde x_r,
\tilde z_1, \ldots, \tilde z_t]
$$
By the same remark, we see that
$D_2 = \det_\kappa(I_\gamma, \alpha, \beta)$ is characterized by
$$
[y_1, \ldots, y_s,
\gamma\tilde x_1, \ldots, \gamma\tilde x_r,
\tilde z_1, \ldots, \tilde z_t]
=
(-1)^{rs} D_2
[x_1, \ldots, x_r,
\gamma\tilde y_1, \ldots, \gamma\tilde y_s,
\tilde z_1, \ldots, \tilde z_t]
$$
Combining the formulas above we see that $D = D_1 D_2$
as desired.
\end{proof}

\end{document}
